% Figures for "Same Text, Different Transitions".
% Every value here is recomputed by artifact/reproduce.sh.

% Figure 1
% The phenomenon: one root, two histories, one rendering, two outcomes.
\newcommand{\FigMechanism}{%
\begin{figure}[t]\centering
\resizebox{\columnwidth}{!}{%
\begin{tikzpicture}[
  font=\small,
  box/.style={draw=black!55,rounded corners=1.5pt,inner sep=3.5pt,
              align=center,fill=white},
  hid/.style={box,draw=BrandA,fill=BrandA!7},
  obs/.style={box,draw=black!70,fill=black!5,thick},
  res/.style={box,draw=BrandB,fill=BrandB!8},
  ar/.style={-{Stealth[length=4pt]},draw=black!55},
]
\node[box] (root) at (0,0) {root state};
\node[hid] (sa) at (2.55,0.72) {$s_1$};
\node[hid] (sb) at (2.55,-0.72) {$s_2$};
\node[obs] (obs) at (4.75,0) {\texttt{$\vdash$ $\uparrow$x + $\uparrow$y = $\uparrow$(x+y)}};
\node[res,anchor=west] (oa) at (6.65,0.72) {accepted};
\node[res,anchor=west] (ob) at (6.65,-0.72) {rejected};

\draw[ar] (root) -- node[above,font=\scriptsize,pos=.45] {$h_1$} (sa);
\draw[ar] (root) -- node[below,font=\scriptsize,pos=.45] {$h_2$} (sb);
\draw[ar] (sa) -- (obs); \draw[ar] (sb) -- (obs);
\draw[ar] (obs.east) to[out=25,in=180] (oa.west);
\draw[ar] (obs.east) to[out=-25,in=180] (ob.west);

\node[font=\scriptsize,BrandA,align=center] at (2.55,1.62)
  {\texttt{ENNReal.ofNNReal x}};
\node[font=\scriptsize,BrandA,align=center] at (2.55,-1.62)
  {\texttt{@WithTop.some x}};
\node[font=\scriptsize,black!60,align=center] at (3.55,-2.28)
  {byte-identical rendering, hence identical model input};
\node[font=\scriptsize,BrandB,align=center,anchor=west] at (6.35,1.62)
  {same tactic};
\end{tikzpicture}}
\caption{Observation aliasing. Two histories from one root reach states that
differ only in elaborated structure the printer collapses; the rendering, and
therefore the entire input supplied to the evaluated state-only generators, is
byte-identical, yet the same tactic is accepted on one and rejected on the
other.}
\label{fig:mechanism}
\end{figure}}

% Figure 2
% Logical evidence ladder with distinct denominators presented separately.
\newcommand{\FigEvidence}{%
\begin{figure}[t]\centering
\resizebox{0.94\columnwidth}{!}{%
\begin{tikzpicture}[
  font=\small,
  ev/.style={draw=black!55,rounded corners=1.5pt,inner sep=4pt,
             text width=6.35cm,align=center,fill=white},
  ar/.style={-{Stealth[length=4pt]},draw=black!55,thick},
]
\node[ev,draw=BrandA,fill=BrandA!7] (e1) at (0,0)
  {\textbf{Existence: do text aliases occur?}\\[-1pt]\scriptsize sampled roots $\rightarrow$ 70 repeat-confirmed classes};
\node[ev,below=0.34cm of e1] (e2)
  {\textbf{Action relevance: do labels differ?}\\[-1pt]\scriptsize confirmed classes $\rightarrow$ 103 contradictory action pairs};
\node[ev,below=0.34cm of e2] (e3)
  {\textbf{Cause and learning: can structure explain the labels?}\\[-1pt]\scriptsize valid interventions; grouped learned acceptance prediction};
\node[ev,draw=BrandB,fill=BrandB!8,below=0.34cm of e3] (e4)
  {\textbf{Consequence: can retention discard a proof?}\\[-1pt]\scriptsize admitted pair--orders; one current-prover case; zero loss with digest};
\draw[ar] (e1) -- (e2); \draw[ar] (e2) -- (e3); \draw[ar] (e3) -- (e4);
\end{tikzpicture}}
\caption{Conceptual evidence spine. Each arrow asks a separate question and
uses the denominator printed beneath it. Population incidence, contradictory
actions, structural learnability, and conditional search loss are estimated
separately.}
\label{fig:evidence}
\end{figure}}

% Figure 3
% Intervention: construction, validity gate, and per-family result.
\newcommand{\FigIntervention}{%
\begin{figure*}[t]\centering
\begin{tikzpicture}[
  font=\small,
  b/.style={draw=black!55,rounded corners=1.5pt,inner sep=4pt,align=center,fill=white},
  gate/.style={b,draw=black!70,fill=black!4},
  ok/.style={b,draw=BrandB,fill=BrandB!8},
  no/.style={b,draw=black!45,fill=black!7},
  ar/.style={-{Stealth[length=4pt]},draw=black!55},
]
% Four stages fit the two-column width. The footnote line states the probing
% step. The two variant boxes
% share a text width so their edges align and the gate arrows stay symmetric.
\node[b] (goal) at (0,0) {one goal\\[-2pt]\scriptsize fixed binders and env};
\node[b,text width=3.9cm] (a)  at (3.6,0.95) {\scriptsize\texttt{show ENNReal.ofNNReal x}};
\node[b,text width=3.9cm] (bb) at (3.6,-0.95) {\scriptsize\texttt{show WithTop.some x}};
\node[gate,text width=3.1cm] (gate) at (7.9,0)
  {\centering valid \emph{iff} both hold\\[1pt]
   \scriptsize renderings byte-identical\\[-2pt]
   \emph{and} digests differ};
\node[ok,text width=3.1cm] (valid) at (11.8,0)
  {\centering\textbf{valid intervention}\\[-1pt]
   \scriptsize observation fixed,\\[-2pt]hidden structure changed};
\draw[ar] (goal) -- (a); \draw[ar] (goal) -- (bb);
\draw[ar] (a.east) -- ($(gate.west)+(0,0.30)$);
\draw[ar] (bb.east) -- ($(gate.west)+(0,-0.30)$);
\draw[ar] (gate) -- (valid);
\node[font=\scriptsize,black!60,align=center] at (5.9,-1.95)
  {a valid pair is probed identically on both variants;
   pairs failing either check are excluded and reported};
\end{tikzpicture}

\vspace{3pt}
\begin{tikzpicture}
\begin{axis}[
  width=0.62\textwidth, height=2.7cm,
  ybar, ymin=0, ymax=17,
  bar width=8pt, enlarge x limits=0.16,
  symbolic x coords={coercion (pinned),coercion (current),Ne vs Not,eta-expansion},
  xtick=data, tick label style={font=\scriptsize},
  ylabel={pairs}, ylabel style={font=\scriptsize},
  axis lines*=left, axis line style={draw=black!45},
  ymajorgrids, grid style={black!12},
  legend style={font=\scriptsize,at={(1.02,1.0)},anchor=north west,draw=black!30},
  legend cell align=left,
  nodes near coords, nodes near coords style={font=\scriptsize},
]
\addplot[fill=black!16,draw=black!40] coordinates {
  (coercion (pinned),15) (coercion (current),8) (Ne vs Not,10) (eta-expansion,10)};
\addlegendentry{constructed}
\addplot[fill=BrandA!55,draw=BrandA!80] coordinates {
  (coercion (pinned),13) (coercion (current),8) (Ne vs Not,0) (eta-expansion,0)};
\addlegendentry{valid}
\addplot[fill=BrandB!75,draw=BrandB!90] coordinates {
  (coercion (pinned),11) (coercion (current),8) (Ne vs Not,0) (eta-expansion,0)};
\addlegendentry{outcome differs}
\end{axis}
\end{tikzpicture}
\caption{Rendering-preserving interventions. \textbf{Top:} the construction.
A single \texttt{show} step replaces the goal with a definitionally equal term,
so only the elaborated representation changes; a pair is used only if it passes
both checks. \textbf{Bottom:} outcomes. The construction succeeds for the
coercion family under the pinned stack and, with acceptance-sensitive probes,
on a current Lean and Mathlib, where all 8 pairs are accepted on one variant
and rejected on the other. The other two families yield zero valid pairs
because the printer renders their two spellings differently.}
\label{fig:intervention}
\end{figure*}}
